\chapter{MEMORIA DESCRIPTIVA}

\section{Generalidades}

La presente memoria descriptiva corresponde al mini proyecto de instalacion electrica interior de una vivienda ubicada en el Jr. Lima N/S, distrito de Capachica, provincia de Puno, departamento de Puno.

El proyecto comprende el diseno preliminar de los circuitos de alumbrado, tomacorrientes, cargas especiales, tablero de distribucion, sistema de proteccion y puesta a tierra, tomando como referencia el Codigo Nacional de Electricidad -- Utilizacion y la Norma EM.010 Instalaciones Electricas Interiores del Reglamento Nacional de Edificaciones.

\section{Objetivo}

El objetivo del proyecto es proponer una instalacion electrica domiciliaria segura, ordenada y justificable tecnicamente para la vivienda seleccionada, mediante:

\begin{itemize}
  \item Levantamiento de ambientes y cargas electricas.
  \item Elaboracion del cuadro de cargas.
  \item Seleccion preliminar de circuitos, conductores y protecciones.
  \item Verificacion basica de corriente y caida de tension.
  \item Preparacion de planos, metrados y especificaciones tecnicas.
\end{itemize}

\section{Base legal y normativa}

El proyecto se desarrolla considerando las normas indicadas en la Tabla~\ref{tab:base-legal}.

\begin{table}[H]
\centering
\caption{Base legal y normativa del proyecto}
\label{tab:base-legal}
\begin{tabularx}{\textwidth}{p{5cm}X}
\toprule
\textbf{Norma} & \textbf{Aplicacion en el proyecto} \\
\midrule
Codigo Nacional de Electricidad -- Utilizacion & Criterios de seguridad, circuitos, conductores, protecciones, tablero, canalizaciones y puesta a tierra. \\
RNE EM.010 Instalaciones Electricas Interiores & Alcance de instalaciones interiores en edificaciones, componentes y documentacion tecnica. \\
Requisitos de la empresa distribuidora local & Punto de suministro, medidor, acometida o condiciones particulares si el docente lo exige. \\
\bottomrule
\end{tabularx}
\end{table}

\textbf{Criterio normativo:} Se toma como base el CNE-U y la Norma EM.010 para instalaciones interiores, conductores, protecciones, tablero, canalizaciones y puesta a tierra.

\section{Descripcion del entorno del proyecto}

\subsection{Ubicacion}

\begin{longtable}{p{5cm}p{9cm}}
\caption{Datos de ubicacion de la vivienda}\label{tab:ubicacion}\\
\toprule
\textbf{Dato} & \textbf{Descripcion} \\
\midrule
\endfirsthead
\toprule
\textbf{Dato} & \textbf{Descripcion} \\
\midrule
\endhead
Direccion o referencia & Jr. Lima N/S \\
Distrito & Capachica \\
Provincia & Puno \\
Departamento & Puno \\
Tipo de zona & Urbana \\
\bottomrule
\end{longtable}

\subsection{Caracteristicas de la vivienda}

\begin{longtable}{p{5cm}p{9cm}}
\caption{Caracteristicas generales de la vivienda}\label{tab:vivienda}\\
\toprule
\textbf{Caracteristica} & \textbf{Valor} \\
\midrule
\endfirsthead
\toprule
\textbf{Caracteristica} & \textbf{Valor} \\
\midrule
\endhead
Tipo de vivienda & Unifamiliar \\
Numero de pisos & 3 pisos \\
Area aproximada & 40 m2 por piso, considerando dimensiones aproximadas de 5 m x 8 m \\
Material predominante & Ladrillo y hormigon \\
Numero de ambientes & Segun levantamiento arquitectonico de la vivienda \\
Tension de suministro & 220 V, 60 Hz, monofasico, salvo verificacion local \\
Tablero general & Tablero general de distribucion empotrado, ubicado en una zona accesible del primer piso, proximo al ingreso principal y al punto de alimentacion \\
\bottomrule
\end{longtable}

\subsection{Ambientes considerados}

\begin{longtable}{c p{4cm} p{4cm} p{5cm}}
\caption{Ambientes considerados en el diseno}\label{tab:ambientes}\\
\toprule
\textbf{Item} & \textbf{Ambiente} & \textbf{Uso} & \textbf{Observacion} \\
\midrule
\endfirsthead
\toprule
\textbf{Item} & \textbf{Ambiente} & \textbf{Uso} & \textbf{Observacion} \\
\midrule
\endhead
1 & Sala & Alumbrado y tomacorrientes & Uso social y tomacorrientes generales \\
2 & Cocina & Alumbrado, tomacorrientes y posibles cargas especiales & Tomacorrientes para pequenos artefactos \\
3 & Dormitorios & Alumbrado y tomacorrientes & Ambientes distribuidos en los tres pisos \\
4 & Bano & Alumbrado y tomacorriente protegido & Priorizar seguridad por humedad \\
5 & Lavanderia & Tomacorriente y carga especial si aplica & Considerar lavadora \\
\bottomrule
\end{longtable}

\section{Caracteristicas tecnicas del proyecto}

\begin{longtable}{p{5cm}p{9cm}}
\caption{Caracteristicas tecnicas preliminares}\label{tab:caracteristicas}\\
\toprule
\textbf{Elemento} & \textbf{Criterio preliminar} \\
\midrule
\endfirsthead
\toprule
\textbf{Elemento} & \textbf{Criterio preliminar} \\
\midrule
\endhead
Sistema electrico & Monofasico 220 V, 60 Hz \\
Tipo de instalacion & Interior domiciliaria empotrada o superficial, segun vivienda \\
Tablero & Tablero general con interruptor principal y circuitos derivados \\
Circuitos & Alumbrado, tomacorrientes generales, cocina/lavanderia y cargas especiales \\
Conductores & Cobre con aislamiento adecuado para uso en edificaciones \\
Canalizacion & Tuberia PVC, conduit u otra permitida segun especificacion \\
Protecciones & Interruptores termomagneticos y diferencial donde corresponda \\
Puesta a tierra & Electrodo, conductor de proteccion y barra de tierra en tablero \\
\bottomrule
\end{longtable}

\section{Impacto ambiental y seguridad}

La instalacion interior propuesta no debe generar impactos ambientales significativos durante su operacion normal. Durante el montaje se deben controlar residuos de cables, tuberias, empaques y elementos cortantes. El diseno debe priorizar la seguridad de las personas mediante proteccion contra sobrecorriente, contactos indirectos, puesta a tierra y ordenamiento del tablero.

\section{Bases para el diseno}

Para el diseno preliminar se consideran los siguientes criterios:

\begin{itemize}
  \item Potencia instalada por ambiente.
  \item Factor de demanda o simultaneidad, segun criterio del curso.
  \item Corriente de diseno por circuito.
  \item Seleccion de conductor por capacidad de corriente.
  \item Verificacion de caida de tension.
  \item Proteccion termomagnetica compatible con conductor y carga.
  \item Puesta a tierra y conductor de proteccion.
\end{itemize}

\section{Cuadro de cargas}

El cuadro de cargas se desarrolla en el Capitulo II. Los valores definitivos deben obtenerse del levantamiento de cargas de la vivienda y de la herramienta HTML ubicada en:

\begin{center}
\texttt{herramientas/calculadora-instalacion-casa.html}
\end{center}

\section{Planos y detalles}

La relacion de planos se desarrolla en el Capitulo VII.

\section{Financiamiento y plazo de ejecucion}

El mini proyecto tiene fines academicos. Para una ejecucion real se debe validar el diseno con un profesional habilitado y con los requisitos de la empresa distribuidora.

\begin{itemize}
  \item Financiamiento: recursos propios para fines academicos.
  \item Plazo estimado: 7 dias de trabajo.
\end{itemize}

